% Section: Milestone Synchronization
\section{Milestone Synchronization}
\label{sec:milestone-synchronization}

% TODO: Expand this section.
%
% Key topics:
% - Definition of milestones in recursive development contexts
% - Milestone as a semantic boundary, not just a time boundary
% - How milestones gate autonomous agent progression
% - Milestone state representation and validation
% - Integration with human-in-the-loop review cycles
% - Failure handling when milestone criteria are not met

Milestone synchronization is the process of coordinating autonomous agent progression
with human-defined semantic checkpoints. In recursive development systems, a
\textit{milestone} represents a verifiable system state that satisfies a predefined
set of quality, correctness, and alignment criteria.

Rather than allowing agents to proceed indefinitely, Reflector-compliant systems
require agents to demonstrate milestone completion before advancing to subsequent
recursive phases. This gating mechanism ensures that errors introduced in early phases
do not silently propagate through the recursive stack.

\subsection{Milestone Definition}

A milestone $M$ is formally defined as a tuple:

\[
M = (S, V, G, \tau)
\]

where:
\begin{itemize}
  \item $S$ is the expected system state at milestone completion
  \item $V$ is the set of validation invariants that must hold
  \item $G$ is the set of governance approvals required
  \item $\tau$ is the maximum elapsed time before escalation
\end{itemize}

\subsection{Synchronization Protocol}

When an agent reaches a milestone boundary, the following synchronization protocol
is executed:

\begin{enumerate}
  \item The agent halts autonomous progression and emits a milestone completion signal.
  \item The recursive auditing system validates the current state against $V$.
  \item If validation passes, the human-in-the-loop governance layer reviews $G$.
  \item Upon approval, the agent resumes execution for the next recursive phase.
  \item If validation fails or approval is withheld, the agent enters a remediation
    cycle governed by the drift recovery protocol (Section~\ref{sec:drift-failure-modes}).
\end{enumerate}
